%-----------------------------------------------------------
% Parametric CAD Framework
%-----------------------------------------------------------

A fully parametric CAD workflow underpins the entire phantom-head model. 
Instead of fixed dimensions, all critical geometric features, tolerances, and mating conditions were defined through named parameters. 
This approach ensures global consistency, minimizes manual editing, and enables rapid regeneration of the model whenever design changes are required.

The parametrization affects all major components of the assembly. 
Examples include the shell thickness in the measuring region, the diameter of the wire-guide channels, the height and keyed profile of the connector bases, and the mating tolerances that govern how components fit together throughout the model.

Table~\ref{tab:parametric_variables} summarizes the principal parameters exposed in the CAD model, along with their default values and functional roles within the design.

\begin{table}[h!]
    \centering
    \caption{Main parametric variables defined in the phantom-head CAD model.}
    \label{tab:parametric_variables}
    \begin{tabular}{l c p{7cm}}
        \toprule
        \textbf{Parameter} & \textbf{Default Value} & \textbf{Description / Purpose} \\
        \midrule

        \texttt{shell\_thickness} & 3.2 mm & Wall thickness of the shell in the measuring region \\

        \bottomrule
    \end{tabular}
\end{table}

Any of these parameters may be modified to accommodate different fabrication processes, material choices, or experimental designs. 
Once updated, the model remains fully consistent and fabrication-ready without the need for manual geometric edits across related pieces. 

\nota{Revisar después: completar la tabla con los nombres reales de las variables, distinguir cuáles son project-wide (en \texttt{project\_params}) y cuáles son locales por archivo. Indicar también la ubicación exacta de cada una dentro del proyecto, quizás en la misma tabla. REREDACTAR TODO, AGREGAR IDEAS, ARCHIVOS, QUE RELACION TIENEN ENTRE ELLOS, ETC.}
